\section{Chaining：把多个模型串起来做单个模型做不到的事}

讲者在最后一段把视角收回到系统设计：foundation model 真正强的地方，不只是“单模型很聪明”，而是可以把多个模型串起来做一件单个模型做不到的事。这就是 chaining。它的基本思想，是把不同模型的长处组合起来，形成一个更可靠、更通用的系统。\footnote{字幕 01:00:25--01:06:30 对 chaining、program generation、retrieval 和 tool use 有非常完整的说明。}

\subsection{为什么要 chaining}

讲者用一个很直观的例子说明：如果 CLIP 认不出一个新类别，但 GPT 能把这个类别描述清楚，那就可以把 GPT 的描述喂给 CLIP 重新做分类。也就是说，text generation 和 visual recognition 可以互相补位。\footnote{字幕 01:01:10--01:02:40 用 Mima、viaduct 等例子说明描述信息可以帮助分类。}

\begin{figure}[H]
\centering
\includegraphics[width=0.95\textwidth]{slides-images/slide-050.jpg}
\caption{chaining 的直觉：一个模型负责描述，另一个模型负责识别或验证}
\end{figure}

\begin{importantbox}{chaining 的本质}
把一个模型当作另一个模型的“工具”或者“中间表示生成器”。它的价值不是替代单模型，而是把多个模块的互补能力组合起来。
\end{importantbox}

\subsection{VisProg：把问题写成程序}

讲者进一步介绍了 visual programming 的思路：让 GPT 先生成一个程序，再由程序去调用对象检测、分割、计数等视觉工具。这样，复杂问题就可以拆成多个步骤，每一步由最擅长的模型完成。例如“图里总共有几个船上的人”，可以先检测每张图中的人，再把检测结果加总。\footnote{字幕 01:02:40--01:05:50 对 VisProg 的程序生成思路做了详细说明。}

\begin{figure}[H]
\centering
\includegraphics[width=0.95\textwidth]{slides-images/slide-007.jpg}
\caption{chaining 也可以理解为把问题转成程序，再用现成工具逐步求解}
\end{figure}

\begin{knowledgebox}{为什么程序化很自然}
很多视觉问题本来就适合分解成子问题：先找对象，再数数量，再做逻辑判断。把这些步骤显式写成程序，往往比让一个模型直接“一步到位”更稳定。
\end{knowledgebox}

\subsection{retrieval 和例子选择}

chaining 不只是调用工具，还包括 retrieval。给模型什么 examples，会显著影响它的行为；如果能够检索最相关的 in-context examples，系统就会比盲目塞一堆例子更好。讲者把这看作“retrieval for prompting”，本质上和传统检索系统是一类问题。\footnote{字幕 01:05:40--01:06:30 对检索式 few-shot prompting 有直接说明。}

\begin{figure}[H]
\centering
\includegraphics[width=0.95\textwidth]{slides-images/slide-014.jpg}
\caption{chaining 的另一种形式：通过 retrieval 选择更合适的 in-context examples}
\end{figure}

\begin{warningbox}{系统成本会迅速上升}
chaining 能力越强，系统也越重：要调用多个模型、多个检索器、多个 verifier，成本和延迟都会上升。所以它是能力增强，也是工程复杂度的增加。
\end{warningbox}

\subsection{hallucination 与 verifier}

讲者最后强调，模型输出并不总可信。尤其在多模态场景里，hallucination 还是常见问题。一个现实的系统往往不会把单模型的输出直接交给用户，而是先经过 verifier 或其他检查模块。换句话说，chain 的价值不只是“接上更多模型”，也在于“把风险分层”。\footnote{字幕 01:05:30--01:06:30 讨论了 hallucination 和 verifier，强调大模型输出通常会再过一层验证。}

\begin{importantbox}{verifier 的角色}
\begin{itemize}
    \item 检查输出和视觉证据是否一致。
    \item 降低 hallucination 被直接暴露给用户的风险。
    \item 让系统从“单点信任”变成“多点审查”。
\end{itemize}
\end{importantbox}

\begin{figure}[H]
\centering
\includegraphics[width=0.95\textwidth]{slides-images/slide-017.jpg}
\caption{系统化使用 foundation models 时，常常需要在输出前增加验证环节}
\end{figure}

